
\documentclass[11pt]{article}
\usepackage[utf8]{inputenc}
\usepackage[T1]{fontenc}
\usepackage{amsmath,amssymb}
\usepackage[margin=1in]{geometry}
\usepackage{enumitem}
\usepackage{hyperref}
\usepackage{titlesec}
\usepackage{xcolor}
\usepackage{fancyhdr}

\pagestyle{fancy}
\fancyhf{}
\lhead{SAT.4D Deep Dive Pronunciation Guide}
\rhead{\thepage}

\titleformat{\section}{\Large\bfseries}{\thesection.}{0.5em}{}
\titleformat{\subsection}{\large\bfseries}{\thesubsection}{0.5em}{}

\title{SAT.4D Deep Dive Pronunciation Guide for Podcast Hosts}
\author{NotebookLM Team}
\date{\today}

\begin{document}

\maketitle

\tableofcontents
\newpage

\section{Overview}

This guide is designed for AI podcast hosts of the NotebookLM Deep Dive series, focusing on SAT.4D modules and their mathematical language. It includes:
\begin{itemize}
  \item Standard phonetic pronunciations (US Midwestern accent)
  \item Natural narration of equations
  \item Suggested phrases to describe key terms in context
  \item Simplification strategies for audio presentation
\end{itemize}

\section{Pronunciation of Common Symbols}

\begin{itemize}[leftmargin=*]
  \item $\alpha$ — ``AL-fuh''
  \item $\beta$ — ``BAY-tuh''
  \item $\gamma$ — ``GAM-uh''
  \item $\delta$ — ``DEL-tuh''
  \item $\epsilon$ — ``EP-suh-lon''
  \item $\zeta$ — ``ZAY-tuh''
  \item $\eta$   — ``AY-tuh''
  \item $\theta$ — ``THAY-tuh''
  \item $\lambda$— ``LAM-duh''
  \item $\mu$    — ``MYOO''
  \item $\nu$    — ``NOO''
  \item $\xi$    — ``KSEE''
  \item $\pi$    — ``PIE''
  \item $\rho$   — ``ROW''
  \item $\sigma$ — ``SIG-muh''
  \item $\tau$   — ``TOW''
  \item $\phi$   — ``FIE''
  \item $\chi$   — ``KAI''
  \item $\psi$   — ``SY''
  \item $\omega$ — ``oh-MEG-uh''
\end{itemize}

\section{Example Equation Narration}

\subsection*{$F = m a$}
Narrate as: ``eff equals em times ay'' or ``force equals mass times acceleration.''

\subsection*{$E = m c^{2}$}
Narrate as: ``ee equals em see squared'' or ``energy equals mass times the speed of light squared.''

\subsection*{$\nabla \cdot \vec{E} = \dfrac{\rho}{\epsilon_{0}}$}
Narrate as: ``the divergence of E equals rho over epsilon naught.''  
Add context: ``This represents Gauss's Law in electromagnetism.''

\section{Narration Tips}

\begin{enumerate}[label=\arabic*.]
  \item Summarize complex formulas in plain language before or after reading them.
  \item Replace mathematical operators with everyday words: ``square root,'' ``over,'' ``times,'' etc.
  \item Skip redundant indices by describing their meaning (e.g.\ dot product, contraction).
  \item Introduce uncommon Greek letters with a short definition or analogy.
  \item Avoid reading code literals; describe their purpose instead.
  \item Break long equations into logical phrases and pause for clarity.
  \item Emphasize physical meaning over raw symbols.
\end{enumerate}

\end{document}
